\section{Absolutely primitive tensor-indecomposable Heisenberg carriers}
\label{sec:heisenberg-absolute-primitive-carrier}

\paragraph{Status and attribution.}
D-TD24--26, reserved question Q2em, concerns a qualified defining-module
case, not every Heisenberg extraction image.  Root proposed applying the
published Guralnick--Saxl theorem to this case and separately developed
the deleted-module trace argument D-TD22--23.
\texttt{/root/g34\_eightfolds} derived the effective generic dimension
bound, exact quotient-exception reduction and full proof below.
The separate D-TD22--23 dependency is accepted.  The author's complete
proof, source-interface and exact-arithmetic checks are recorded
separately; root and independent whole-proof reviews of this section pass.

\paragraph{Actual image and field hypotheses.}
Use the actual sphere-cover setup of D-TD05--07: a faithful irreducible
finite coefficient image $J\subset\operatorname{GL}(V)$, an actual
connected cover $C\to\mathbf P^1$, and a stable whole factorization
$V=W\otimes R$.  The entire external partner $R$ is retained with
arbitrary dimension and action.  Put $e=\dim V$, and suppose
\[
 \sum_i\operatorname{rank}(1-T_i)=2e+6,\qquad
 P=\operatorname{im}(J\to\operatorname{PGL}(W)),\qquad
 Q\triangleleft P,\quad Q\simeq\mathbf F_p^d,\quad d=2m\ge2.
 \tag{GHC.1}
\]
The projective $Q$-action is irreducible with nondegenerate scalar
commutator form.  D-TD05 gives
\[
 a=\dim W=p^{d/2},\qquad
 H=P/Q\ \hookrightarrow\operatorname{Sp}(Q),\qquad
 \frac{|\operatorname{tr}A|}{a}
 \in\{0,p^{-\operatorname{rank}(1-h)/2}\}
 \tag{GHC.2}
\]
for a unitary representative $A$ of a projective element inducing
$h\in H$.  If $h=1$ and the projective element is nonidentity,
the trace is zero.  In particular $H$ is the actual faithful
linear action on $Q$, not an abstract composition factor.

Whenever the Guralnick--Saxl hypotheses are imposed below, they are
imposed on
\[
                         \overline Q=Q\otimes_{\mathbf F_p}
                                      \overline{\mathbf F}_p:
 \qquad \overline Q\text{ is irreducible, primitive and
 tensor-indecomposable as an }H\text{-module}.
 \tag{GHC.3}
\]
Primitivity and tensor indecomposability only over $\mathbf F_p$
are not substituted for (GHC.3).  Tensor indecomposability has the
published meaning of no nontrivial factorwise preserved tensor
decomposition; the source explicitly allows a group permuting an
equal-factor tensor structure.  No hypothesis of semisimplicity in
the defining characteristic is silently added.

\subsection*{D-TD24. The generic Guralnick--Saxl branch has dimension below 576}

\begin{proposition}[D-TD24]
In (GHC.1)--(GHC.2), suppose every nonscalar $h\in H$ satisfies
\[
                  \operatorname{rank}_{\mathbf F_p}(1-h)
                                      >\frac{\sqrt d}{2}.
 \tag{GHC.4}
\]
Then $d<576$.  In particular this conclusion holds in the
nonexceptional branch of Guralnick--Saxl Theorem~7.1 under (GHC.3).
\end{proposition}

\begin{proof}
Suppose $d\ge576$.  Rank is unchanged by scalar extension.
Every nonscalar nonidentity $h$ has integer rank at least thirteen
by (GHC.4), since $\sqrt{576}/2=12$.  A nonidentity scalar $h$
has rank $d$, so the same bound holds.  A nonidentity element of
$P$ with $h=1$ has zero trace by D-TD05.  Therefore every
nonidentity projective element has a unitary representative with
normalized absolute trace at most
\[
                 c=p^{-13/2}\le2^{-13/2}<\frac1{90},
                 \qquad 2^{13}=8192>8100=90^2.
 \tag{GHC.5}
\]

Let $Z=C/\ker(J\to P)$.  It is an actual connected $P$-Galois
cover of the sphere.  Its genus is at least two: $m=d/2\ge288$,
so the genus-zero and genus-one elementary-abelian rank exclusions
in D-TD06 apply to its faithfully acting subgroup $Q$.
Remove only identity images from its inertia tuple, and let
$o_1,\ldots,o_k$ be the exact orders of the remaining elements
of $P$.  For a projective element of order $o_i$, choose a unitary
representative with $o_i$th power one.  All its powers numbered
$1,\ldots,o_i-1$ are nonscalar, so Fourier inversion and (GHC.5)
give all-target normalized support at least
$(1-c)(1-1/o_i)$.  Spectral convolution preserves this lower
bound after tensoring by the entire $R$, regardless of reciprocal
scalar choices.  Identity-projected original entries contribute
nonnegative actual rank and can be omitted in this lower bound.

Actual Riemann--Hurwitz and the elementary positive orbifold
bound in D-TD06 give $\sum_i(1-1/o_i)\ge85/42$.  Thus
\[
       2+\frac6e\ge(1-c)\frac{85}{42},\qquad
       e(1-85c)\le252,\qquad e<4536.
 \tag{GHC.6}
\]
The last inequality follows from $1-85c>1/18$.  On the other
hand $e\ge a=p^{d/2}\ge2^{288}>4536$, a contradiction.
This proves the first assertion.

Here is the exact published theorem interface for the last assertion.
Over an algebraically closed field, Guralnick--Saxl define
$\nu_H(\overline Q)$ to be the minimum rank of $1-\lambda h$
over $h\in H$ and scalars $\lambda$ for which $\lambda h\ne I$.
For dimension $d\ge6$ and (GHC.3), their Theorem~7.1 says either
\[
             \nu_H(\overline Q)>\max\{2,\sqrt d/2\},
 \tag{GHC.7}
\]
or one of its listed exceptions holds.  The first branch implies
(GHC.4) by taking $\lambda=1$.  The precise input is the theorem
and Remark~3, pp.~548--549, with its complete supplied proof
pp.~549--554.\cite{BI:GS}
The root in (GHC.7) is $\sqrt d$ divided by two, not
$\sqrt{d/2}$.  The source proof's explicit deduction
$d\le\alpha(H)^2/4$ from $\nu\alpha\ge d$ confirms this
reading of its extracted formula.  No fresh reconstruction of
the upstream classification or character tables is asserted.
\end{proof}

\subsection*{D-TD25. An actual large alternating quotient leaves only the deleted module}

\begin{theorem}[D-TD25]
Assume (GHC.1)--(GHC.3), and suppose that the actual conjugation
image has a surjection
\[
             \psi:H\longrightarrow T,\qquad
             T\in\{A_n,S_n\},\qquad n\ge10^6.
 \tag{GHC.8}
\]
Then necessarily $p=2$, $\psi$ is an isomorphism, and $Q$ is
the fully deleted permutation module of that actual $A_n$ or
$S_n$, of dimension $n-1$ or $n-2$.
\end{theorem}

\begin{proof}
First $p<7744$.  Otherwise D-TD07 applied to the actual whole
factor $P$ gives $P=E[p]\rtimes C_s$ with
$s\in\{2,3,4,6\}$, which is solvable.  Its quotient $H$
cannot then surject onto the non-solvable $T$.  This uses the
actual $P$-cover and the entire external partner, exactly as
in D-TD19, and imposes no projected-mark lower bound.

If $d<576$, its evenness gives $d\le574$.  Counting all matrices
over $\mathbf F_p$ supplies the deliberately crude bound
\[
 |H|\le|\operatorname{GL}_d(p)|<p^{d^2}
       <2^{13\cdot574^2}=2^{4283188}
       <2^{8999999}<\frac{n!}{2}\le|T|.
 \tag{GHC.9}
\]
For the last strict inequality, the last $500000$ factors of
$1000000!$ are each greater than $2^{18}$, and factorial is
increasing for $n\ge1000000$.  This contradicts surjectivity.
No order formula for a classical group is needed here.
We may therefore assume $d\ge576$.

D-TD24 rules out the first branch (GHC.7).  Theorem~7.1 and its
Remark~3 leave only the natural classical case or the deleted
alternating/symmetric case: every other exception has dimension
at most ten.  We spell out the normal-subgroup information used
from that theorem, rather than treating an arbitrary subgroup of
a classical group as a full classical core.  Its proof on p.~549
first reduces to $S=F^*(H/Z(H))$ simple.  The subsequent natural
classical and alternating cases identify this simple group and
its natural action.  The normal quasisimple component can also
be recovered directly: let $K$ be the preimage of $S$ in $H$
and take the stable term $E=K^{(\infty)}$ of its derived series.
Then $E$ is perfect, normal in $H$, maps onto the perfect $S$,
and has central kernel.  Hence it is quasisimple.

Its restriction on $\overline Q$ is semisimple: the nonzero
$E$-socle is $H$-invariant, hence the whole irreducible module.
Its isotypic summands are permuted by $H$, so absolute
primitivity leaves only one.  Evaluation gives an irreducible
$E$-space tensor its multiplicity space, with both full factor
algebras preserved by $H$.  The first space has dimension
greater than one because $E$ is perfect and acts faithfully.
Tensor indecomposability therefore makes the multiplicity one.
Thus $E$ acts absolutely irreducibly, and its projective simple
quotient is the full named classical group or $A_c$, respectively.

For completeness, the needed outer quotient follows directly
from this component.  Schur's lemma makes $C_H(E)=Z(H)$.
An automorphism of the perfect $E$ which induces the identity
on $S=E/Z(E)$ differs from the identity by a central character,
and is therefore trivial.  Consequently
\[
       H/(EZ(H))\ \hookrightarrow\operatorname{Out}(S),
       \qquad H/E\text{ is solvable}.
 \tag{GHC.10}
\]
The last assertion uses the published Schreier theorem, with
exact locator Babai--Luks--Seress, Theorem~2.3 p.~1314;
its classification dependence is retained.\cite{HD:BabaiLuksSeress1997}
The central part of $H/E$ is abelian.  Thus no assertion about
all subgroups or all sections of $\operatorname{Aut}(S)$ is used.

The image $\psi(E)$ is perfect and normal in $T$.  If it were
trivial, $T$ would be a quotient of the solvable $H/E$.
Hence $\psi(E)=A_n$, and quasisimplicity identifies $S$ with
$A_n$.  The natural classical case is impossible.  The precise
classification input is the complete repetition list in
Li--Xia, Section~2.2.1, numbered author p.~13: the only
alternating/classical coincidences are
$A_5=\operatorname{PSL}_2(4)=\operatorname{PSL}_2(5)$,
$A_6=\operatorname{PSL}_2(9)$ and
$A_8=\operatorname{PSL}_4(2)$.\cite{HES:LiXia2022}
None has $n\ge10^6$.  This is an explicit classification input,
not a blanket inference from the labels of the two families.

It remains to examine the deleted-module case.  In that case
$E$ is the actual ordinary $A_c$ on its fully deleted module.
Indeed the deleted action is an ordinary representation of
$A_c$, and the derived group of its scalar extension is that
perfect $A_c$ itself; a different central extension cannot be
manufactured by scalar choices.  Since $\psi(E)\simeq A_n$
and $E\simeq A_c$, equality of their factorial orders gives
$c=n$.

Here is the form argument, including the defining characteristic.
Over $F=\overline{\mathbf F}_p$ let
\[
 M=F^c,\qquad M_0=\{(v_i):\sum_i v_i=0\},\qquad
 D=M_0/(M_0\cap F(1,\ldots,1)).
 \tag{GHC.11}
\]
The standard dot product on $M$ induces a nondegenerate symmetric
form on $D$, since its radical on $M_0$ is precisely the line
being divided out, when that line lies in $M_0$.
The module $D$ is absolutely irreducible for $A_c$ here.  An
elementary verification is useful: lift a nonzero invariant
submodule to $M_0$, take a vector which is not constant, and
subtract a suitable three-cycle translate.  This gives a nonzero
vector supported on at most three coordinates.  Choose one
nonzero coordinate and three distinct zero coordinates; subtract
its translate under a double transposition pairing the first
two and the other two.  The result is a nonzero multiple of
$e_i-e_j$.  The alternating orbit of this vector spans $M_0$.
There are enough zero coordinates since $c=n\ge10^6$.

The given nondegenerate alternating form on $Q$, after scalar
extension, is also $E$-invariant.  Irreducibility and Schur's
lemma make it a scalar multiple of the nondegenerate symmetric
dot-product form.  In odd characteristic that is impossible.
Hence $p=2$.  (In characteristic two the induced dot product
is alternating on $M_0$.)

Finally, $\psi|_E$ is injective.  For $h\in\ker\psi$ its
commutators with $E$ lie both in $E$ and in $\ker\psi$, so
$h$ centralizes $E$.  Its matrix is scalar on $\overline Q$,
and, since its entries lie in $\mathbf F_2$, that scalar is one.
The actual action of $H$ is faithful; thus $h=1$ and
\(\psi\) is an isomorphism.

We check the exact prime-field module and form, not just their
scalar extensions.  The identification of the alternating core
through $\psi$ differs from the standard one by an automorphism
of $A_n$.  Since $n\ge10^6$, every such automorphism is induced
by a permutation of the $n$ letters.  Relabeling therefore makes
the two core modules isomorphic over $\overline{\mathbf F}_2$.
Their intertwiner equations are linear equations over
$\mathbf F_2$; scalar extension preserves the dimension of
their solution space.  That space has dimension one by
irreducibility, so it already contains a nonzero
$\mathbf F_2$-intertwiner, which is an isomorphism.
After this identification, the matrix of any $h\in H$ and
the natural deleted matrix of $\psi(h)$ induce the same
conjugation on the core.  Their ratio centralizes that core
and is a scalar in $\mathbf F_2^*=\{1\}$.  Thus their
matrices agree on the entire $H$.  The two invariant
nondegenerate forms likewise differ by an element of
$\mathbf F_2^*$, hence coincide.  We have proved exactly
\[
                   H\simeq T,\qquad p=2,\qquad
                   Q\text{ is its fully deleted symplectic module}.
 \tag{GHC.12}
\]
This is an actual quotient isomorphism and actual module, not
merely an alternating composition factor of a symplectic group.
\end{proof}

\subsection*{D-TD26. Qualified elimination of the carrier}

\begin{corollary}[D-TD26, using D-TD22--23]
There is no actual sphere source satisfying (GHC.1)--(GHC.3)
and the actual surjection (GHC.8).  Thus an alternating quotient
carrier of natural degree $n\ge10^6$ arising in D-TD19 cannot
have an absolutely irreducible defining module which is primitive
and tensor-indecomposable after scalar extension to
$\overline{\mathbf F}_p$.
\end{corollary}

\begin{proof}
D-TD25 forces precisely the actual characteristic-two deleted
module treated in D-TD22--23.  That separate result excludes its
stable Heisenberg factor at $n\ge10^6$ with the same arbitrary
entire external partner and exact original inertia budget.
All of its quotient and field hypotheses have therefore been
verified, so it gives the contradiction.
\end{proof}

\paragraph{Retained boundary.}
The conclusion is conditional on the stated absolute primitivity
and tensor indecomposability, not on a purported implication
from their $\mathbf F_p$ analogues.  Nonabsolutely irreducible,
absolutely imprimitive and absolutely tensor-decomposable defining
modules remain outside this result.  No bound on their exponent
$m$, no independent original subgroup lift, and no equality of
the Heisenberg image with the full symplectic group is inferred.
